% !TeX root = ../main.tex
\begin{center}
    \textbf{\LARGE ДОМАШНЯЯ ПОДГОТОВКА} \\[0.5cm]
    {\large «Основные параметры фильтров высокочастотного диапазона»}
\end{center}

\subsection*{1. Расчет мощности и коэффициента затухания кабеля}

1. Переведем мощность генератора в дБм относительно опорного уровня $P_0 = 1\text{ мВт}$:
$$P_{\text{ген }} = 10 \lg \left(\frac{P_{\text{ген [мВт]}}}{1\text{ мВт}}\right) = 10 \lg(10) = 10\text{ дБм}$$

2. Мощность сигнала на выходе аттенюатора с ослаблением $A_{\text{атт}} = 10\text{ дБ}$:
$$P_{\text{атт [дБм]}} = P_{\text{ген [дБм]}} - A_{\text{атт}} = 10 - 10 = 0\text{ дБм}$$

Переведем полученное значение в линейные единицы мощности (Вт):
$$P_{\text{атт [мВт]}} = 10^{\frac{P_{\text{атт [дБм]}}}{10}} = 10^{0} = 1\text{ мВт}$$
$$P_{\text{атт [Вт]}} = 1 \cdot 10^{-3}\text{ Вт} = 0{,}001\text{ Вт}$$

3. Коэффициент затухания кабеля $L_{\text{каб}}$ определяется как разность уровней мощности на его входе (выходе аттенюатора) и выходе (входе анализатора спектра):
$$L_{\text{каб [дБ]}} = P_{\text{атт [дБм]}} - P_{\text{ан [дБм]}} = 0 - (-13) = 13\text{ дБ}$$


